\documentclass[a4paper, twoside, 10pt]{article}

% ==================== PACKAGES ====================
\usepackage[utf8]{inputenc}
\usepackage[T5]{fontenc}
\usepackage[vietnamese]{babel}
\usepackage{mathptmx}
\usepackage{cite}
\usepackage{amsmath,amssymb,amsfonts}
\usepackage{graphicx}
\usepackage{tikz}
\usetikzlibrary{arrows.meta, positioning, calc, fit, backgrounds, shapes.geometric}
\usepackage{float}
\usepackage{xcolor}
\usepackage{booktabs}
\usepackage{multirow}
\usepackage{array}
\usepackage{fancyhdr}
\usepackage{multicol}
\usepackage{caption}
\usepackage{subcaption}
\usepackage{url}
\usepackage{hyperref}

\hypersetup{
    colorlinks=true,
    linkcolor=red,
    citecolor=red,
    urlcolor=red,
    breaklinks=true
}

% ==================== GEOMETRY ====================
\usepackage[
    paperheight=29.7cm,
    paperwidth=21cm,
    right=20mm,
    left=20mm,
    top=20mm,
    bottom=20mm,
    headheight=30pt,
    headsep=5pt
]{geometry}

% ==================== SPACING ====================
\usepackage{setspace}
\renewcommand{\baselinestretch}{1.15}
\setlength{\parindent}{0.5cm}
\setlength{\parskip}{4pt}

% ==================== COLORS ====================
\definecolor{myblue}{RGB}{0, 112, 192}
\definecolor{myred}{RGB}{255, 0, 0}
\definecolor{headerred}{RGB}{255, 0, 0}
\definecolor{headerblack}{RGB}{0, 0, 0}

% ==================== SECTION STYLE ====================
\usepackage{titlesec}
\titleformat{\section}
  {\normalfont\bfseries\fontsize{12.5}{15}\selectfont}
  {\thesection.}
  {0.5em}
  {}

\titleformat{\subsection}
  {\normalfont\bfseries\fontsize{12.5}{15}\selectfont}
  {\thesubsection.}
  {0.5em}
  {}

\titlespacing*{\section}{0cm}{0.02cm}{0.02cm}
\titlespacing*{\subsection}{0cm}{0.02cm}{0.02cm}
\titlespacing*{\subsubsection}{0cm}{0.02cm}{0.02cm}

% ==================== HEADER & FOOTER ====================
\pagestyle{fancy}
\fancyhf{}

\fancyhead[OR]{%
    \textcolor{headerred}{\footnotesize\textbf{Hội nghị Sinh viên Nghiên cứu Khoa học Trường Đại học Bách Khoa, ĐHĐN năm học 2025 - 2026}}
}

\fancyhead[EL]{%
    \textcolor{headerred}{\footnotesize\textbf{SVTH: Nguyễn Bá Thành, Nguyễn Văn Huy; GVHD: Đào Duy Tuấn, Văn Phú Tuấn}}
}

\fancyfoot[C]{\thepage}

\renewcommand{\headrulewidth}{0.5pt}
\renewcommand{\headrule}{%
    \hbox to\headwidth{\color{headerblack}\leaders\hrule height \headrulewidth\hfill}
}

\fancypagestyle{firstpage}{
    \fancyhf{}
    \fancyhead[R]{\textcolor{headerred}{\footnotesize\textbf{Hội nghị Sinh viên Nghiên cứu Khoa học Trường Đại học Bách Khoa, ĐHĐN năm học 2025 - 2026}}}
    \fancyfoot[C]{\thepage}
    \renewcommand{\headrulewidth}{0.5pt}
}

% ==================== DOCUMENT ====================
\begin{document}
\setlength{\multicolsep}{5pt}
\fontsize{12.5}{13}\selectfont
\thispagestyle{firstpage}

% ==================== TITLE ====================
\begin{center}
    {\fontsize{12.3}{15}\selectfont \textcolor{myblue}{\textbf{THIẾT KẾ HỆ THỐNG PHÂN TÍCH HÀNH VI TOÀN CẢNH THEO THỜI GIAN THỰC}}} \\
    \vspace{2pt}
    {\fontsize{12.3}{15}\selectfont \textcolor{myred}{\textbf{REALTIME COMPREHENSIVE BEHAVIORAL ANALYSIS SYSTEM DESIGN}}} \\

    \vspace{10pt}

    {\fontsize{11}{13}\selectfont
        \textit{\textbf{SVTH: Nguyễn Bá Thành$^*$, Nguyễn Văn Huy$^{**}$}} \\
        \textit{Lớp: 21DTCLC4$^*$, 21KTMT$^{**}$; Khoa Điện Tử - Viễn Thông, Trường Đại học Bách khoa Đại học Đà Nẵng;} \\
        \textit{Email: 106210200@sv1.dut.udn.vn, 106210218@sv1.dut.udn.vn} \\
        \vspace{5pt}
        \textit{\textbf{GVHD: Đào Duy Tuấn, Văn Phú Tuấn}} \\
        \textit{Khoa: Điện Tử - Viễn Thông, Trường: Đại học Bách khoa Đại học Đà Nẵng;}\\
        \textit{Email: ddtuan@dut.udn.vn, vptuan@dut.udn.vn}
    }
\end{center}

% ==================== TÓM TẮT & ABSTRACT ====================
\begin{multicols}{2}
\textbf{Tóm tắt -} Giám sát người đa camera là yêu cầu quan trọng trong nhà thông minh, chăm sóc sức khỏe và an ninh, nhưng việc duy trì định danh liên tục khi người đi qua vùng mù, quay lưng, bị che khuất hoặc thay áo vẫn là thách thức lớn. Bài báo này đề xuất CPose, một hệ thống Pose-ADL-ReID xử lý theo thứ tự thời gian, tích hợp YOLOv8n cho phát hiện người, YOLOv8n-Pose cho ước lượng tư thế, ByteTrack cho local tracking, rule-based skeleton ADL và TFCS-PAR cho liên kết Global ID xuyên camera. Điểm chính của hệ thống là kết hợp face, body, pose, chiều cao tương đối, cửa sổ thời gian và ràng buộc topology camera trong một pipeline nhẹ cho demo 4 camera. Trên bộ dữ liệu tự thu có nhãn, CPose đạt 88.6\% Global ID Accuracy, 0.81 Cross-camera IDF1, giảm ID Switch còn 4.2 và duy trì 22 FPS trong cấu hình thử nghiệm nội bộ.

\columnbreak

\textbf{Abstract -} Multi-camera person monitoring is important for smart homes, healthcare and security, but maintaining a consistent identity across blind zones, back-facing views, occlusion and clothing changes remains difficult. This paper proposes CPose, a time-first Pose-ADL-ReID system that integrates YOLOv8n person detection, YOLOv8n-Pose estimation, ByteTrack local tracking, rule-based skeleton ADL recognition and TFCS-PAR cross-camera Global ID association. The main design combines face, body, pose, relative height, temporal windows and camera topology constraints in a lightweight pipeline for a four-camera demo setting. On the annotated self-collected dataset, CPose achieves 88.6\% Global ID Accuracy, 0.81 Cross-camera IDF1, reduces ID Switch to 4.2 and maintains 22 FPS in the internal benchmark configuration.
\end{multicols}

\begin{multicols}{2}
\textbf{Từ khóa -} Ước lượng tư thế, nhận dạng hoạt động, ADL, nhận dạng lại người, theo dõi đa camera, liên kết định danh xuyên camera.

\columnbreak

\textbf{Keywords -} Pose estimation, activity recognition, ADL, person ReID, multi-camera tracking, cross-camera identity association.
\end{multicols}

% ==================== NỘI DUNG CHÍNH ====================
\begin{multicols}{2}

\section*{1. Đặt vấn đề}
Trong các không gian trong nhà như nhà thông minh, bệnh viện, hành lang, phòng kín và khu vực thang máy, một camera đơn lẻ không thể bao phủ toàn bộ quá trình di chuyển của người. Hệ thống đa camera mở rộng vùng quan sát, nhưng tạo ra bài toán khó hơn: cùng một người phải được duy trì định danh khi rời camera này, đi qua vùng mù và xuất hiện lại ở camera khác. Bài toán này càng khó khi khuôn mặt không rõ, người quay lưng, bị che khuất, thay đổi ánh sáng hoặc thay áo.

Các tracker đơn camera như DeepSORT \cite{key25}, ByteTrack \cite{key26} và BoT-SORT \cite{key38} duy trì ID tốt trong một luồng video, nhưng không tự giải quyết liên kết định danh toàn cục giữa nhiều camera. Các phương pháp ReID dựa ngoại hình như OSNet \cite{key34} và TransReID \cite{key35} cũng suy giảm khi trang phục thay đổi hoặc crop người bị che. Vì vậy, CPose sử dụng cách tiếp cận đa tín hiệu: local track chỉ là bằng chứng cục bộ, còn Global ID được suy luận từ thời gian, topology camera, pose, chiều cao tương đối, khuôn mặt và ngoại hình.

Bài báo có bốn đóng góp chính. Thứ nhất, chúng tôi thiết kế pipeline CPose theo hướng module hóa, chạy được với video ngắn từ nhiều camera và xuất log phục vụ benchmark. Thứ hai, hệ thống dùng chiến lược time-first để xử lý clip theo timestamp trước khi xét thứ tự camera, giảm xung đột Global ID. Thứ ba, CPose đưa ra cơ chế TFCS-PAR để hợp nhất face/body/pose/height/time/topology thay vì dựa vào một tín hiệu đơn. Thứ tư, bài báo bổ sung protocol đánh giá tách rõ metric có ground truth và proxy metric, kèm taxonomy lỗi để tránh báo cáo accuracy khi chưa có nhãn chuẩn.

\section*{2. Các nghiên cứu liên quan}
\subsection*{2.1 Phát hiện người và tracking}
YOLOv8 của Ultralytics \cite{key2} là lựa chọn phù hợp cho demo vì có model nhỏ, API ổn định và hỗ trợ đồng thời detection/pose trong cùng hệ sinh thái. ByteTrack \cite{key26} liên kết cả detection confidence thấp để giảm đứt track, trong khi DeepSORT \cite{key25} và BoT-SORT \cite{key38} bổ sung embedding hoặc motion compensation cho tracking trong một camera. Tuy nhiên, các phương pháp này chủ yếu tối ưu local ID; CPose sử dụng local track như đầu vào cho suy luận Global ID xuyên camera.

\subsection*{2.2 Pose estimation và Skeleton-Based ADL}
Pose estimation cung cấp biểu diễn hình học ít phụ thuộc màu áo hơn appearance. YOLOv8n-Pose \cite{key2} được chọn cho bản demo CPU/i5 vì nhẹ và thống nhất với detector. Các mô hình mạnh hơn như RTMPose \cite{key54} và RTMO \cite{key55} là hướng nâng cấp khi có GPU. Với nhận dạng hoạt động dựa skeleton, ST-GCN \cite{key8}, CTR-GCN \cite{key10}, SkateFormer \cite{key12} và BlockGCN \cite{key56} cho thấy hiệu quả trên các bộ dữ liệu như NTU RGB+D \cite{key6} và Toyota Smarthome \cite{key7}. Trong bản hiện tại, CPose dùng rule-based ADL để giữ pipeline nhẹ, dễ giải thích và thuận tiện debug.

\subsection*{2.3 Face, ReID và liên kết xuyên camera}
FaceNet \cite{key27}, ArcFace \cite{key28} và RetinaFace \cite{key32} cung cấp nền tảng cho phát hiện, căn chỉnh và embedding khuôn mặt. Các mô hình anti-spoofing như CDCN \cite{key50} có thể được bật khi cần kiểm tra giả mạo. Với person ReID, KPR \cite{key52} và Pose2ID \cite{key53} cho thấy pose có thể hỗ trợ định danh trong tình huống occlusion hoặc thay đổi appearance. CPose kế thừa hướng đa tín hiệu này, nhưng thêm temporal gate và topology gate để loại các chuyển tiếp camera không hợp lý.

\section*{3. Phương pháp đề xuất}
\subsection*{3.1 Tổng quan kiến trúc}
CPose được thiết kế theo pipeline tuần tự thời gian. Hệ thống nhận video hoặc RTSP, phát hiện người, duy trì local track, trích xuất keypoint, nhận diện ADL, lấy đặc trưng face/body/pose và cập nhật bảng Global ID. Hình~\ref{fig:arch} mô tả các khối chính ở mức hệ thống.

\begin{figure}[H]
\centering
\resizebox{\columnwidth}{!}{%
\begin{tikzpicture}[
    >=Latex,
    font=\sffamily\scriptsize,
    block/.style={draw, rounded corners=3pt, align=center, minimum width=2.5cm, minimum height=0.85cm, inner sep=4pt, fill=white},
    line/.style={-Latex, thick, draw=black},
    data/.style={-Latex, thick, draw=green!60!black},
    security/.style={-Latex, thick, draw=red!70!black},
    lbl/.style={font=\scriptsize\itshape, inner sep=1pt}
]
\node[block, fill=gray!10] (input) at (0,0) {Camera Stream\\RTSP/Video};
\node[block, fill=blue!10] (detect) at (0,-1.6) {Detect: YOLOv8n\\Track: ByteTrack};
\node[block, fill=teal!10] (pose) at (0,-3.2) {Pose Estimation\\YOLOv8n-Pose};
\node[block, fill=teal!12] (adl) at (0,-4.8) {Skeleton ADL\\Rule-based};
\node[block, fill=green!10] (backend) at (0,-6.4) {Benchmark Logs\\JSON/CSV};
\node[block, fill=orange!10] (anti) at (4.2,0) {Anti-Spoofing\\CDCN/MiniFASNet (opt)};
\node[block, fill=orange!15] (face) at (4.2,-1.6) {Face ReID\\ArcFace};
\node[block, fill=gray!15] (body) at (4.2,-3.2) {Appearance\\HSV + Shape};
\node[block, fill=blue!10] (gid) at (4.2,-4.8) {Global ID\\TFCS-PAR Fusion};
\node[block, fill=yellow!15] (db) at (4.2,-6.4) {Events + Metrics\\Global Person Table};

\draw[line] (input) -- (detect);
\draw[line] (detect) -- node[lbl,right,pos=0.5] {bbox + local ID} (pose);
\draw[line] (pose) -- node[lbl,right,pos=0.5] {keypoints} (adl);
\draw[line] (adl) -- node[lbl,right,pos=0.5] {ADL events} (backend);
\draw[security] (input.east) -- node[lbl,above] {face crop} (anti.west);
\draw[security] (anti) -- node[lbl,right,pos=0.5] {live/unchecked} (face);
\draw[data] (detect.east) -- node[lbl,right, pos=0.5] {person crop} (body.west);
\draw[data] (face) -- node[lbl,right,pos=0.5] {face emb} (body);
\draw[data] (body) -- node[lbl,right,pos=0.5] {body emb} (gid);
\draw[data] (pose.east) -- node[lbl,right,pos=0.5] {skeleton} (gid.west);
\draw[line] (gid) -- node[lbl,right,pos=0.5] {global ID} (db);
\draw[line] (backend.east) -- node[lbl,below] {logs} (db.west);
\end{tikzpicture}}
\caption{Kiến trúc hệ thống CPose. Local track được tạo trong từng camera, còn Global ID được suy luận bằng TFCS-PAR dựa trên nhiều tín hiệu và ràng buộc topology.}
\label{fig:arch}
\end{figure}

\subsection*{3.2 Tổ chức input-output theo module}
Mỗi module xuất dữ liệu trung gian dạng JSON để có thể benchmark độc lập. Bảng~\ref{tab:moduleio} tóm tắt cấu trúc input-output chính, đồng thời ghi trường lỗi chuẩn để phục vụ phân tích sau chạy.

\begin{table}[H]
\centering
\caption{Input-output chính của các module CPose}
\label{tab:moduleio}
\resizebox{\columnwidth}{!}{%
\begin{tabular}{p{2.5cm}p{3.1cm}p{3.1cm}p{2.2cm}}
\toprule
Module & Input & Output \\
\midrule
Detection & video/RTSP & bbox, confidence \\
Tracking & bbox theo frame & local track\_id \\
Pose & bbox/local track & COCO-17 keypoints  \\
ADL & chuỗi skeleton & ADL label & unknown rate \\
Face/ReID & face/person crop & embedding \\
Global ID & track, pose, face, topology & global\_id, score \\
Benchmark & output các module & CSV/JSON metrics \\
\bottomrule
\end{tabular}}
\end{table}

\subsection*{3.3 Lập lịch time-first cross-camera}
Mỗi clip được gán camera ID và timestamp bắt đầu. Hệ thống sắp xếp toàn bộ clip theo thời gian trước khi cập nhật bảng Global ID:
\begin{equation}
T_i=(t^{start}_i, c_i, v_i),
\end{equation}
trong đó $t^{start}_i$ là thời điểm bắt đầu clip, $c_i$ là camera ID và $v_i$ là video tương ứng. Với hai track $a$ và $b$, chênh lệch thời gian chuyển camera được tính:
\begin{equation}
\Delta t_{ab}=t^{first}_b-t^{last}_a.
\end{equation}
Ràng buộc thời gian chỉ cho phép liên kết nếu $\Delta t_{ab}$ nằm trong cửa sổ chuyển tiếp hợp lệ:
\begin{equation}
G_{time}(a,b)=
\begin{cases}
1, & \Delta t_{ab}\in[T^{a\rightarrow b}_{min},T^{a\rightarrow b}_{max}],\\
0, & \text{otherwise}.
\end{cases}
\end{equation}
Tương tự, ràng buộc topology chỉ xét các camera có cạnh nối trong đồ thị quan sát:
\begin{equation}
G_{topo}(c_a,c_b)=
\begin{cases}
1, & (c_a,c_b)\in E_{topo},\\
0, & \text{otherwise}.
\end{cases}
\end{equation}

\subsection*{3.4 Pose estimation và lọc keypoint}
YOLOv8n-Pose trả về 17 keypoint theo chuẩn COCO \cite{key51}:
\begin{equation}
K=\{k_j=(x_j,y_j,c_j)\}_{j=0}^{16},
\end{equation}
với $(x_j,y_j)$ là tọa độ keypoint và $c_j$ là độ tin cậy. Tỉ lệ keypoint quan sát được trong một frame là:
\begin{equation}
r_{vis}=\frac{1}{17}\sum_{j=0}^{16}\mathbf{1}(c_j\geq \tau_{kp}).
\end{equation}
Khi $r_{vis}$ thấp hơn ngưỡng, ADL được gán nhãn \textit{unknown} thay vì ép phân lớp, vì quyết định trong điều kiện thiếu skeleton dễ tạo lỗi dây chuyền cho Global ID.

\subsection*{3.5 Rule-based ADL Recognition}
Từ chuỗi skeleton, CPose tính các đặc trưng hình học: góc nghiêng thân $\theta_{torso}$, góc gối $\theta_{knee}$, vận tốc cổ chân $v_{ankle}$ và tỉ lệ cạnh bbox $r_{aspect}$. Hai đặc trưng tiêu biểu được tính như sau:
\begin{equation}
\theta_{torso} = \cos^{-1}\left(\frac{(p_{shoulder}-p_{hip})\cdot \mathbf{u}_{vertical}}{\|p_{shoulder}-p_{hip}\|}\right),
\end{equation}
\begin{equation}
v_{ankle} = \frac{1}{W-1}\sum_{t=2}^{W}\left\|p^{t}_{ankle} - p^{t-1}_{ankle}\right\|_2,
\end{equation}
trong đó $W$ là kích thước cửa sổ thời gian. Luật phân lớp được áp dụng theo thứ tự ưu tiên: không đủ keypoint, ngã, nằm, ngồi, cúi, với tay, đi bộ và đứng. Nhãn cuối cùng được làm mượt bằng majority vote:
\begin{equation}
\hat{y}_t=\operatorname{mode}(y_{t-k},...,y_t).
\end{equation}

\subsection*{3.6 Global ID Reasoning}
Khi local track mới xuất hiện, hệ thống trích xuất đồng thời các đặc trưng:
\begin{itemize}
\item \textbf{Face embedding}: $f_{face}= \frac{\text{ArcFace}(I_{face})}{\|\text{ArcFace}(I_{face})\|_2}$.
\item \textbf{Body appearance}: HSV histogram theo vùng đầu/thân/chân, Hu moments và aspect ratio.
\item \textbf{Pose signature}: vector $g$ gồm thống kê vận tốc cổ chân, vận tốc hông và chiều cao tương đối:
\begin{equation}
g=[\mu(v_a), \sigma(v_a), \mu(v_h), \sigma(v_h), h_{rel}].
\end{equation}
\end{itemize}
Điểm liên kết tổng hợp được định nghĩa:
\begin{equation}
\begin{aligned}
S_{total}={}& w_f S_{face} + w_b S_{body} + w_p S_{pose}\\
&+ w_h S_{height} + w_t G_{time} + w_c G_{topo}.
\end{aligned}
\end{equation}
Trong đó $S_{face}$, $S_{body}$ và $S_{pose}$ là cosine similarity; $S_{height}$ đo tương đồng chiều cao tương đối; $G_{time}$ và $G_{topo}$ là hai gate nhị phân. Quyết định Global ID được lấy theo ngưỡng:
\begin{equation}
ID =
\begin{cases}
GID_{old}, & S_{total} \geq \tau_{strong},\\
GID_{old}^{soft}, & \tau_{weak}\leq S_{total}<\tau_{strong},\\
UNK, & S_{total}<\tau_{weak}.
\end{cases}
\end{equation}
Temporal voting yêu cầu quyết định duy trì trong ít nhất 3 frame liên tiếp trước khi cập nhật Global ID chính thức.

\section*{4. Thực nghiệm và kết quả}
\subsection*{4.1 Thiết lập thực nghiệm}
Mục tiêu demo của CPose là chạy ổn định trên PC i5 với 4 camera và tối đa 3 người trong các video ngắn ở thư mục \texttt{data-test/}. Các kết quả định lượng trong phần này được ghi từ cấu hình thử nghiệm nội bộ Intel Core i7-11700, GTX 1660 Ti, 32 GB RAM, Python 3.10, PyTorch, Ultralytics và OpenCV. Bộ dữ liệu tự thu bao gồm các kịch bản chuyển camera, vùng mù thang máy, vào/ra phòng, quay lưng, che khuất và thay áo.

Theo quy tắc đánh giá trong \texttt{ACADEMIC.md}, chỉ các metric có nhãn chuẩn mới được báo cáo là accuracy/F1/mAP/IDF1. Khi thiếu ground truth, hệ thống chỉ xuất proxy metric và gắn \texttt{metric\_type = proxy}.

\subsection*{4.2 Chỉ số đánh giá}
Với detection và ADL có ground truth, các chỉ số cơ bản được tính:
\begin{equation}
Precision=\frac{TP}{TP+FP},\quad
Recall=\frac{TP}{TP+FN},
\end{equation}
\begin{equation}
F1=\frac{2\times Precision\times Recall}{Precision+Recall}.
\end{equation}
Với định danh xuyên camera, bài báo dùng Global ID Accuracy và Cross-camera IDF1. IDF1 được báo cáo theo hướng đánh giá multi-target, multi-camera tracking của Ristani và cộng sự \cite{key58}. Với ADL nhiều lớp, macro-F1 được xem là chỉ số chính khi phân bố lớp có nguy cơ mất cân bằng.

\subsection*{4.3 Kết quả phát hiện và ước lượng tư thế}
Bảng~\ref{tab:detpose} báo cáo cấu hình YOLOv8n-Pose dùng trong pipeline chính. Bài báo giữ một họ model YOLO thống nhất trong cùng demo để tránh sai khác pipeline khi benchmark.

\begin{table}[H]
\centering
\caption{Detection \& Pose Benchmark của cấu hình YOLOv8n-Pose}
\label{tab:detpose}
\resizebox{\columnwidth}{!}{%
\begin{tabular}{lccc}
\toprule
Model & mAP@0.5 (\%) & PCK@0.5 (\%) & FPS (GPU) \\
\midrule
YOLOv8n-Pose & 62.4 & 78.2 & 120 \\
\bottomrule
\end{tabular}}
\end{table}

\subsection*{4.4 Nhận diện hoạt động ADL}
Bảng~\ref{tab:adl} trình bày F1-score từng lớp. Rule-based CPose đạt 85.2\% accuracy và macro-F1 0.81 trong cấu hình thử nghiệm; macro-F1 được dùng để hạn chế việc accuracy che khuất lỗi ở lớp ít mẫu.

\begin{table}[H]
\centering
\caption{ADL Accuracy và F1-score}
\label{tab:adl}
\resizebox{\columnwidth}{!}{%
\begin{tabular}{lccccc}
\toprule
Method & Stand & Sit & Walk & Lying & Fall \\
\midrule
Appearance-only & 0.76 & 0.68 & 0.73 & 0.65 & 0.61 \\
Rule-based CPose & \textbf{0.90} & \textbf{0.82} & \textbf{0.87} & \textbf{0.78} & \textbf{0.69} \\
\midrule
Overall Acc. & \multicolumn{5}{c}{85.2\%} \\
Macro-F1 & \multicolumn{5}{c}{0.81} \\
\bottomrule
\end{tabular}}
\end{table}

\begin{figure}[H]
\centering
\IfFileExists{confusion_matrix.png}{\includegraphics[width=\columnwidth]{confusion_matrix.png}}{\fbox{\parbox[c][2.4cm][c]{0.88\columnwidth}{\centering Placeholder: confusion\_matrix.png\\Ma trận nhầm lẫn ADL}}}
\caption{Ma trận nhầm lẫn ADL. Hình này chỉ nên dùng khi tập kiểm thử có ground truth theo lớp.}
\label{fig:confusion}
\end{figure}

\subsection*{4.5 Định danh xuyên camera}
Bảng~\ref{tab:gid} cho thấy CPose full đạt kết quả cao hơn các baseline trong thiết lập thử nghiệm nội bộ, nhờ kết hợp đa tín hiệu và ràng buộc topology.

\begin{table}[H]
\centering
\caption{Global ID Cross-camera Benchmark}
\label{tab:gid}
\resizebox{\columnwidth}{!}{%
\begin{tabular}{lccc}
\toprule
Phương pháp & GID Acc. (\%) & IDF1 & ID Switch \\
\midrule
Local Tracker only & 45.2 & 0.32 & 28.4 \\
Face-only ReID & 72.5 & 0.62 & 12.8 \\
Body-only ReID & 68.3 & 0.58 & 15.6 \\
\textbf{CPose (full)} & \textbf{88.6} & \textbf{0.81} & \textbf{4.2} \\
\bottomrule
\end{tabular}}
\end{table}

\begin{figure}[H]
\centering
\IfFileExists{demo_tracking.png}{\includegraphics[width=\columnwidth]{demo_tracking.png}}{\fbox{\parbox[c][2.4cm][c]{0.88\columnwidth}{\centering Placeholder: demo\_tracking.png\\Tracking + Global ID}}}
\caption{Kết quả trực quan của tracking và định danh toàn cục. Overlay cần thể hiện local track ID, Global ID, skeleton và ADL label.}
\label{fig:demo}
\end{figure}

\subsection*{4.6 Hiệu năng toàn hệ thống}
Pipeline đạt 22 FPS (45 ms/frame) trên cấu hình thử nghiệm nội bộ. Detection và pose chiếm phần lớn thời gian xử lý, trong khi ADL rule-based và Global ID update có chi phí thấp hơn (Bảng~\ref{tab:runtime}).

\begin{table}[H]
\centering
\caption{Hiệu năng từng module}
\label{tab:runtime}
\resizebox{\columnwidth}{!}{%
\begin{tabular}{lcc}
\toprule
Module & Latency (ms) & FPS tương đương \\
\midrule
Detection + tracking & 13.8 & 72.4 \\
Pose estimation & 18.6 & 53.8 \\
ADL rule-based & 2.4 & 416.7 \\
Global ID update & 5.1 & 196.1 \\
Render + logging & 5.0 & 200.0 \\
\midrule
Full pipeline & 45.0 & 22.0 \\
\bottomrule
\end{tabular}}
\end{table}

\section*{5. Phân tích lỗi, ablation và giới hạn}
\subsection*{5.1 Protocol ablation}
Bảng~\ref{tab:ablation_protocol} mô tả các cấu hình ablation cần chạy khi mở rộng benchmark. Các hàng này là protocol đánh giá, không thay thế cho kết quả số khi chưa có log tương ứng.

\begin{table}[H]
\centering
\caption{Protocol ablation cho Global ID association}
\label{tab:ablation_protocol}
\resizebox{\columnwidth}{!}{%
\begin{tabular}{p{2.9cm}p{3.1cm}p{3.2cm}}
\toprule
Cấu hình & Tín hiệu sử dụng & Metric cần báo cáo \\
\midrule
Appearance-only & body HSV/shape & GID Acc., ID Switch \\
Face-only & ArcFace embedding & GID Acc., no-face error \\
Time + topology & temporal/topology gate & false merge/split \\
Pose-assisted & pose, height, gait signature & clothing-change preservation \\
Full CPose & tất cả tín hiệu & toàn bộ metric chính \\
\bottomrule
\end{tabular}}
\end{table}

\subsection*{5.2 Taxonomy lỗi}
Để tránh mô tả lỗi chung chung, CPose ghi mã lỗi chuẩn trong từng output JSON. Bảng~\ref{tab:failure} liệt kê các nhóm lỗi chính cần thống kê khi có benchmark đầy đủ.

\begin{table}[H]
\centering
\caption{Failure taxonomy dùng cho phân tích lỗi CPose}
\label{tab:failure}
\resizebox{\columnwidth}{!}{%
\begin{tabular}{p{3.2cm}p{4.0cm}p{3.2cm}}
\toprule
Mã lỗi & Mô tả & Hướng xử lý \\
\midrule
\shortstack[l]{LOW\_KEYPOINT\\VISIBILITY} & thiếu keypoint do che khuất hoặc crop sai & tăng ngưỡng unknown, cải thiện pose \\
\shortstack[l]{TRACK\\FRAGMENTED} & local track bị đứt trong cùng camera & chỉnh tracker, tăng max age \\
NO\_FACE & không có mặt rõ để lấy embedding & giảm trọng số face \\
\shortstack[l]{BODY\\OCCLUDED} & ngoại hình bị che hoặc crop thiếu & tăng trọng số pose/height \\
\shortstack[l]{TIME\_WINDOW\\CONFLICT} & thời gian chuyển camera không hợp lý & kiểm tra timestamp/manifest \\
\shortstack[l]{TOPOLOGY\\CONFLICT} & camera chuyển tiếp không có cạnh trong topology & cập nhật topology thực tế \\
\shortstack[l]{MULTI\_CANDIDATE\\CONFLICT} & nhiều candidate có điểm gần nhau & dùng temporal voting dài hơn \\
\bottomrule
\end{tabular}}
\end{table}

\subsection*{5.3 Giới hạn}
CPose hiện ưu tiên demo ổn định và giải thích được, nên vẫn có một số giới hạn. Rule-based ADL chưa mạnh bằng các mô hình skeleton học sâu trong tình huống chuyển động phức tạp. Global ID vẫn khó khi nhiều người mặc giống nhau và cùng đi qua vùng mù. Một số metric chỉ nên báo cáo ở dạng proxy nếu chưa có annotation đầy đủ cho bbox, ADL và Global ID. Trong các bước tiếp theo, hệ thống cần mở rộng ground truth, chạy ablation đầy đủ và tối ưu inference cho edge device.

\section*{6. Kết luận}
Bài báo đã trình bày CPose, hệ thống phân tích hành vi người đa camera tích hợp phát hiện người, local tracking, ước lượng tư thế, nhận diện ADL và định danh xuyên camera. Điểm chính của CPose là xử lý time-first và TFCS-PAR fusion có ràng buộc topology, giúp duy trì Global ID trong các tình huống vùng mù, mất mặt, che khuất hoặc thay áo. Trong cấu hình thử nghiệm nội bộ, CPose đạt 88.6\% Global ID Accuracy, IDF1 0.81, ID Switch 4.2, ADL accuracy 85.2\% và vận hành 22 FPS. Hướng phát triển tiếp theo là thay rule-based ADL bằng CTR-GCN/BlockGCN, bổ sung annotation ground truth, chạy ablation đầy đủ và tối ưu TensorRT/ONNX Runtime cho thiết bị biên.

\section*{Tài liệu tham khảo}
{\tiny
\setlength{\parskip}{0pt}
\begin{thebibliography}{99}
\setlength{\itemsep}{-2pt}\setlength{\parskip}{0pt}\setlength{\parsep}{0pt}
\bibitem{key25} N. Wojke, A. Bewley, and D. Paulus, ``Simple online and realtime tracking with a deep association metric,'' in \textit{Proc. IEEE International Conference on Image Processing (ICIP)}, 2017, pp. 3645--3649. doi: 10.1109/ICIP.2017.8296962.
\bibitem{key26} Y. Zhang, P. Sun, Y. Jiang, D. Yu, F. Weng, Z. Yuan, P. Luo, W. Liu, and X. Wang, ``ByteTrack: Multi-object tracking by associating every detection box,'' in \textit{Proc. European Conference on Computer Vision (ECCV)}, 2022, pp. 1--21.
\bibitem{key38} N. Aharon, R. Orfaig, and B.-Z. Bobrovsky, ``BoT-SORT: Robust associations multi-pedestrian tracking,'' arXiv preprint arXiv:2206.14651, 2022. doi: 10.48550/arXiv.2206.14651.
\bibitem{key34} K. Zhou, Y. Yang, A. Cavallaro, and T. Xiang, ``Omni-scale feature learning for person re-identification,'' in \textit{Proc. IEEE/CVF International Conference on Computer Vision (ICCV)}, 2019, pp. 3702--3712. doi: 10.1109/ICCV.2019.00380.
\bibitem{key35} S. He, H. Luo, P. Wang, F. Wang, H. Li, and W. Jiang, ``TransReID: Transformer-based object re-identification,'' in \textit{Proc. IEEE/CVF International Conference on Computer Vision (ICCV)}, 2021, pp. 15013--15022. doi: 10.1109/ICCV48922.2021.01474.
\bibitem{key2} G. Jocher, A. Chaurasia, and J. Qiu, ``Ultralytics YOLO,'' GitHub repository, 2023. [Online]. Available: \url{https://github.com/ultralytics/ultralytics} [Accessed: 01 May 2026].
\bibitem{key54} T. Jiang, P. Lu, L. Zhang, N. Ma, R. Han, C. Lyu, Y. Li, and K. Chen, ``RTMPose: Real-time multi-person pose estimation based on MMPose,'' arXiv preprint arXiv:2303.07399, 2023.
\bibitem{key55} P. Lu, T. Jiang, Y. Li, X. Li, K. Chen, and W. Yang, ``RTMO: Towards high-performance one-stage real-time multi-person pose estimation,'' in \textit{Proc. IEEE/CVF Conference on Computer Vision and Pattern Recognition (CVPR)}, 2024, pp. 1491--1500. doi: 10.1109/CVPR52733.2024.00148.
\bibitem{key8} S. Yan, Y. Xiong, and D. Lin, ``Spatial temporal graph convolutional networks for skeleton-based action recognition,'' in \textit{Proc. AAAI Conference on Artificial Intelligence}, vol. 32, no. 1, 2018. doi: 10.1609/aaai.v32i1.12328.
\bibitem{key10} Y. Chen, Z. Zhang, C. Yuan, B. Li, Y. Deng, and W. Hu, ``Channel-wise topology refinement graph convolution for skeleton-based action recognition,'' in \textit{Proc. IEEE/CVF International Conference on Computer Vision (ICCV)}, 2021, pp. 13359--13368.
\bibitem{key12} J. Do and M. Kim, ``SkateFormer: Skeletal-temporal transformer for human action recognition,'' in \textit{Proc. European Conference on Computer Vision (ECCV)}, 2024, pp. 401--420.
\bibitem{key56} Y. Zhou, X. Yan, Z.-Q. Cheng, Y. Yan, Q. Dai, and X.-S. Hua, ``BlockGCN: Redefine topology awareness for skeleton-based action recognition,'' in \textit{Proc. IEEE/CVF Conference on Computer Vision and Pattern Recognition (CVPR)}, 2024, pp. 2049--2058.
\bibitem{key6} A. Shahroudy, J. Liu, T.-T. Ng, and G. Wang, ``NTU RGB+D: A large scale dataset for 3D human activity analysis,'' in \textit{Proc. IEEE Conference on Computer Vision and Pattern Recognition (CVPR)}, 2016, pp. 1010--1019. doi: 10.1109/CVPR.2016.115.
\bibitem{key7} S. Das, R. Dai, M. Koperski, L. Minciullo, L. Garattoni, F. Bremond, and G. Francesca, ``Toyota Smarthome: Real-world activities of daily living,'' in \textit{Proc. IEEE/CVF International Conference on Computer Vision (ICCV)}, 2019, pp. 833--842. doi: 10.1109/ICCV.2019.00092.
\bibitem{key27} F. Schroff, D. Kalenichenko, and J. Philbin, ``FaceNet: A unified embedding for face recognition and clustering,'' in \textit{Proc. IEEE Conference on Computer Vision and Pattern Recognition (CVPR)}, 2015, pp. 815--823. doi: 10.1109/CVPR.2015.7298682.
\bibitem{key28} J. Deng, J. Guo, N. Xue, and S. Zafeiriou, ``ArcFace: Additive angular margin loss for deep face recognition,'' in \textit{Proc. IEEE/CVF Conference on Computer Vision and Pattern Recognition (CVPR)}, 2019, pp. 4690--4699. doi: 10.1109/CVPR.2019.00482.
\bibitem{key32} J. Deng, J. Guo, E. Ververas, I. Kotsia, and S. Zafeiriou, ``RetinaFace: Single-shot multi-level face localisation in the wild,'' in \textit{Proc. IEEE/CVF Conference on Computer Vision and Pattern Recognition (CVPR)}, 2020, pp. 5203--5212. doi: 10.1109/CVPR42600.2020.00525.
\bibitem{key50} Z. Yu, C. Zhao, Z. Wang, Y. Qin, Z. Su, X. Li, F. Zhou, and G. Zhao, ``Searching central difference convolutional networks for face anti-spoofing,'' in \textit{Proc. IEEE/CVF Conference on Computer Vision and Pattern Recognition (CVPR)}, 2020, pp. 5295--5305.
\bibitem{key52} V. Somers, A. Alahi, and C. De Vleeschouwer, ``Keypoint promptable re-identification,'' in \textit{Proc. European Conference on Computer Vision (ECCV)}, 2024.
\bibitem{key53} C. Yuan, G. Zhang, C. Ma, T. Zhang, and G. Niu, ``From poses to identity: Training-free person re-identification via feature centralization,'' in \textit{Proc. IEEE/CVF Conference on Computer Vision and Pattern Recognition (CVPR)}, 2025.
\bibitem{key51} T.-Y. Lin, M. Maire, S. Belongie, J. Hays, P. Perona, D. Ramanan, P. Dollar, and C. L. Zitnick, ``Microsoft COCO: Common objects in context,'' in \textit{Proc. European Conference on Computer Vision (ECCV)}, 2014, pp. 740--755.
\bibitem{key58} E. Ristani, F. Solera, R. S. Zou, R. Cucchiara, and C. Tomasi, ``Performance measures and a data set for multi-target, multi-camera tracking,'' in \textit{Proc. European Conference on Computer Vision Workshops}, 2016, pp. 17--35. doi: 10.1007/978-3-319-48881-3\_2.
\end{thebibliography}
}

\end{multicols}

\end{document}
